\chapter{Transfer Matrices and Scattering}
\label{app:transfer}

This appendix develops the one-dimensional transfer-matrix formalism
that appears in Chapters~\ref{ch:pt-symmetric}, \ref{ch:poles-zeros},
and~\ref{ch:bbc-rebuilt}. The results are exact for multilayer
dielectric stacks with complex permittivity and provide the
computational backbone for many of the non-Hermitian scattering
calculations in this book.


%% ================================================================
\section{Single-layer transfer matrix}
\label{sec:single-layer}

Consider a homogeneous layer of thickness $d$, refractive index
$n = n_r + i n_i$, at normal incidence. The electric field at
position $x$ is a superposition of forward and backward waves:
\begin{equation}
 E(x) = A\,e^{ink_0 x} + B\,e^{-ink_0 x},
 \label{eq:field-ab}
\end{equation}
where $k_0 = \omega/c$. The transfer matrix $M$ relates the field
amplitudes at the left face ($x = 0$) to those at the right face
($x = d$):
\begin{equation}
 \begin{pmatrix} A \\ B \end{pmatrix}_{\!x=d}
 = M\,
 \begin{pmatrix} A \\ B \end{pmatrix}_{\!x=0},
 \qquad
 M = \begin{pmatrix}
 e^{ink_0 d} & 0 \\
 0 & e^{-ink_0 d}
 \end{pmatrix}.
 \label{eq:propagation-matrix}
\end{equation}
At an interface between media with indices $n_1$ and $n_2$, the
interface matrix is
\begin{equation}
 I_{12} = \frac{1}{2n_2}\begin{pmatrix}
 n_2 + n_1 & n_2 - n_1 \\
 n_2 - n_1 & n_2 + n_1
 \end{pmatrix}.
 \label{eq:interface-matrix}
\end{equation}
The total transfer matrix for a multilayer stack is the ordered product
of propagation and interface matrices:
\begin{equation}
 \Tmat = I_{N,\mathrm{out}} \cdot M_N \cdots I_{23} \cdot M_2
 \cdot I_{12} \cdot M_1.
 \label{eq:total-transfer}
\end{equation}


%% ================================================================
\section{Scattering matrix from the transfer matrix}
\label{sec:transfer-to-scatter}

The $2 \times 2$ scattering matrix $\Smat$ relates the outgoing
amplitudes to the incoming amplitudes:
\begin{equation}
 \begin{pmatrix} B_L \\ A_R \end{pmatrix}
 = \Smat
 \begin{pmatrix} A_L \\ B_R \end{pmatrix}
 = \begin{pmatrix}
 r_L & t' \\
 t & r_R
 \end{pmatrix}
 \begin{pmatrix} A_L \\ B_R \end{pmatrix},
 \label{eq:smat-def}
\end{equation}
where $r_L$, $r_R$ are the reflection coefficients from the left and
right, and $t$, $t'$ are the transmission coefficients. The conversion
from the transfer matrix $\Tmat = \bigl(\begin{smallmatrix}
T_{11} & T_{12} \\ T_{21} & T_{22} \end{smallmatrix}\bigr)$ is
\begin{equation}
 t = \frac{1}{T_{22}}, \qquad
 t' = \frac{\det\Tmat}{T_{22}}, \qquad
 r_L = -\frac{T_{21}}{T_{22}}, \qquad
 r_R = \frac{T_{12}}{T_{22}}.
 \label{eq:transfer-to-scatter}
\end{equation}

For a reciprocal medium (symmetric $\epsr$ tensor), $t = t'$ and
$\det\Tmat = 1$. For a non-reciprocal medium, $t \neq t'$ in general.


%% ================================================================
\section{$\PT$-symmetric transfer matrices}
\label{sec:pt-transfer-app}

A one-dimensional structure with $\PT$ symmetry satisfies
$n(x) = n^*(-x)$. For a multilayer consisting of a loss layer
(index $n_0 + i\kappa$, thickness $d$) followed by a gain layer
(index $n_0 - i\kappa$, thickness $d$), the unit-cell transfer
matrix is
\begin{equation}
 \Tmat_{\mathrm{cell}} = M_{\mathrm{gain}} \cdot I \cdot
 M_{\mathrm{loss}}.
 \label{eq:pt-unit-cell}
\end{equation}
The $\PT$ condition imposes the constraint
\begin{equation}
 \sigma_x\,\Tmat_{\mathrm{cell}}^*\,\sigma_x = \Tmat_{\mathrm{cell}},
 \qquad \sigma_x = \begin{pmatrix} 0 & 1 \\ 1 & 0 \end{pmatrix},
 \label{eq:pt-transfer-constraint}
\end{equation}
which ensures that $T_{11} = T_{22}^*$ and $T_{12} = T_{21}^*$. For
the scattering matrix, this gives
\begin{equation}
 |t|^2 - |r_L|\,|r_R| = 1
 \quad\text{(exact $\PT$ relation)},
 \label{eq:pt-unitarity}
\end{equation}
which replaces the Hermitian unitarity condition $|t|^2 + |r|^2 = 1$.


\subsection{$\PT$ phase transition from transfer-matrix eigenvalues}
\label{ssec:pt-phases}

The eigenvalues of $\Tmat_{\mathrm{cell}}$ are
$\lambda_\pm = \tfrac{1}{2}\mathrm{tr}(\Tmat_{\mathrm{cell}})
\pm \sqrt{[\tfrac{1}{2}\mathrm{tr}(\Tmat_{\mathrm{cell}})]^2 - 1}$.
Three regimes arise:
\begin{enumerate}
 \item \textbf{$\PT$-symmetric phase}
 ($|\tfrac{1}{2}\mathrm{tr}\,\Tmat| \le 1$):
 $\lambda_\pm$ are unimodular; the stack is a Bragg grating with
 real band structure.
 \item \textbf{$\PT$-broken phase}
 ($|\tfrac{1}{2}\mathrm{tr}\,\Tmat| > 1$):
 $|\lambda_+| > 1 > |\lambda_-|$; one eigenvalue grows exponentially
 with the number of cells (laser-like amplification), while the
 other decays (CPA-like absorption).
 \item \textbf{Exceptional point}
 ($|\tfrac{1}{2}\mathrm{tr}\,\Tmat| = 1$):
 $\lambda_+ = \lambda_-$; the transfer matrix has a Jordan block
 of size 2. This is the $\PT$-symmetry breaking threshold.
\end{enumerate}


%% ================================================================
\section{Poles and zeros from the transfer matrix}
\label{sec:poles-zeros-transfer}

The poles and zeros of the scattering matrix can be located directly
from the transfer matrix:
\begin{itemize}
 \item A \textbf{pole} of $\Smat(\omega)$ (a lasing mode) occurs when
 $T_{22}(\omega) = 0$, since $t = 1/T_{22}$.
 \item A \textbf{zero} of $\det\Smat(\omega)$ (a CPA mode) occurs when
 $T_{11}(\omega)=0$, since with the convention above
 $\det\Smat=T_{11}/T_{22}$. For reciprocal media $\det\Tmat=1$, but
 this does not prevent $\det\Smat$ from vanishing because the scattering
 determinant is not the same object as the transfer determinant.
\end{itemize}

The CPA-laser point is the frequency at which both $T_{22} = 0$ and
$T_{11} = 0$ simultaneously. For a $\PT$-symmetric stack at real
frequency, the constraint $T_{11}=T_{22}^*$ makes the pole and zero
conditions coincide; in practice this simultaneous condition is isolated
in the $(\omega,\kappa)$ parameter space.


%% ================================================================
\section{Bloch waves in periodic stacks}
\label{sec:bloch-transfer}

For an infinite periodic stack with unit-cell transfer matrix
$\Tmat_{\mathrm{cell}}$, the Bloch condition is
\begin{equation}
 \Tmat_{\mathrm{cell}}\,
 \begin{pmatrix} A \\ B \end{pmatrix}
 = e^{ika}\,
 \begin{pmatrix} A \\ B \end{pmatrix},
 \label{eq:bloch-transfer}
\end{equation}
so the Bloch phase is $\beta = e^{ik a}$ and the dispersion relation is
\begin{equation}
 \cos(ka) = \tfrac{1}{2}\,\mathrm{tr}(\Tmat_{\mathrm{cell}}).
 \label{eq:dispersion-transfer}
\end{equation}
For a Hermitian stack, $\mathrm{tr}(\Tmat_{\mathrm{cell}})$ is real,
and band gaps correspond to $|\cos(ka)| > 1$. For a non-Hermitian
stack, $\mathrm{tr}(\Tmat_{\mathrm{cell}})$ is complex, and the
dispersion defines a curve in the complex $\omega$ plane---the complex
band structure of Chapter~\ref{ch:complex-bands}.

\subsection{GBZ from the transfer matrix}
\label{ssec:gbz-transfer}

For a finite open stack, the correct Brillouin zone is the GBZ
(Section~\ref{sec:gbz}). The GBZ condition $|\beta_n| = |\beta_{n+1}|$
translates to the condition that the two eigenvalues of
$\Tmat_{\mathrm{cell}}$ have equal modulus:
$|\lambda_+| = |\lambda_-|$. Since $\det\Tmat_{\mathrm{cell}} =
\lambda_+\lambda_- = 1$ for reciprocal media, this gives
$|\lambda_\pm| = 1$---the GBZ coincides with the ordinary BZ, and there
is no skin effect. For non-reciprocal media,
$|\det\Tmat_{\mathrm{cell}}| \neq 1$, and the GBZ deforms away from
the unit circle, producing the skin effect.

This transfer-matrix derivation of the GBZ is the most direct route
for one-dimensional photonic structures and connects the abstract
non-Bloch theory of Chapter~\ref{ch:skin-effect} to concrete
multilayer calculations.


%% ================================================================
\section{Numerical implementation}
\label{sec:transfer-numerical}

\subsection{Avoiding numerical overflow}
\label{ssec:overflow}

For thick stacks with gain ($n_i < 0$), the propagation matrix
elements $e^{\pm i n k_0 d}$ can overflow or underflow. The standard
remedy is to factor the transfer matrix after each layer, extracting
the largest eigenvalue modulus and accumulating the logarithm of the
total transmission. This ``scattering-matrix propagation algorithm''
is numerically stable and avoids the exponential overflow that plagues
naive transfer-matrix products.

\subsection{Complex-frequency evaluation}
\label{ssec:complex-freq}

For QNM and pole-zero calculations, the transfer matrix must be
evaluated at complex frequencies $\omega = \omega_r + i\omega_i$.
The refractive index $n(\omega)$ is then evaluated by analytic
continuation of the dispersion model (Lorentz, Drude, etc.), and the
propagation phases $n k_0 d$ become fully complex. No additional
numerical modification is needed beyond ensuring that the branch cuts
of the square root in $n(\omega) = \sqrt{\epsr(\omega)}$ are
consistently chosen.


%% ================================================================
\section{From transfer matrix to scattering matrix}
\label{sec:t-to-s}

The transfer matrix and scattering matrix encode the same physics but
organise it differently. Converting between them is a routine but
essential operation in photonic multilayer calculations.

\subsection{Conversion formulas}
\label{ssec:ts-conversion}

For a single interface or slab described by the $2\times 2$ transfer
matrix
\begin{equation}
 \Tmat = \begin{pmatrix} T_{11} & T_{12} \\ T_{21} & T_{22} \end{pmatrix},
 \label{eq:transfer-generic}
\end{equation}
relating the field amplitudes on the left to those on the right,
$\bigl(\begin{smallmatrix} A_R \\ B_R \end{smallmatrix}\bigr)
= \Tmat\,\bigl(\begin{smallmatrix} A_L \\ B_L \end{smallmatrix}\bigr)$,
the scattering matrix
$\bigl(\begin{smallmatrix} B_L \\ A_R \end{smallmatrix}\bigr)
= \Smat\,\bigl(\begin{smallmatrix} A_L \\ B_R \end{smallmatrix}\bigr)$
is obtained by rearranging:
\begin{equation}
 \Smat = \frac{1}{T_{22}}
 \begin{pmatrix}
 -T_{21} & 1 \\
 T_{11}T_{22} - T_{12}T_{21} & T_{12}
 \end{pmatrix}
 = \frac{1}{T_{22}}
 \begin{pmatrix}
 -T_{21} & 1 \\
 \det\Tmat & T_{12}
 \end{pmatrix}.
 \label{eq:t-to-s}
\end{equation}
The elements have direct physical meaning:
$t_L = 1/T_{22}$ is the left-to-right transmission amplitude,
$r_L = -T_{21}/T_{22}$ is the left reflection, and
$r_R = T_{12}/T_{22}$ is the right reflection. For a reciprocal
medium, $\det\Tmat = 1$ and the left-to-right and right-to-left
transmission amplitudes are equal: $t_L = t_R = 1/T_{22}$.

\subsection{Example: Fabry--P\'erot resonator}
\label{ssec:fp-from-transfer}

Consider a dielectric slab of refractive index $n$ and thickness $d$
in air. The total transfer matrix is the product of three matrices:
the entrance interface, the propagation through the slab, and the exit
interface. Using the Fresnel interface matrices and the propagation
matrix from Section~\ref{sec:single-layer}:
\begin{equation}
 \Tmat_{\mathrm{FP}} = \Tmat_{\mathrm{exit}}\,\Tmat_{\mathrm{prop}}\,\Tmat_{\mathrm{entrance}}.
 \label{eq:fp-transfer-product}
\end{equation}
The entrance interface matrix (from air to the slab, normal incidence)
is
\begin{equation}
 \Tmat_{\mathrm{entrance}} = \frac{1}{2n}
 \begin{pmatrix} n+1 & n-1 \\ n-1 & n+1 \end{pmatrix},
 \label{eq:entrance-interface}
\end{equation}
the propagation matrix through the slab is
\begin{equation}
 \Tmat_{\mathrm{prop}} = \begin{pmatrix}
 e^{ink_0 d} & 0 \\ 0 & e^{-ink_0 d}
 \end{pmatrix},
 \label{eq:fp-propagation-matrix}
\end{equation}
where $k_0 = \omega/c$, and the exit interface matrix (from the slab
to air) is
\begin{equation}
 \Tmat_{\mathrm{exit}} = \frac{1}{2}
 \begin{pmatrix} 1+1/n & 1-1/n \\ 1-1/n & 1+1/n \end{pmatrix}.
 \label{eq:exit-interface}
\end{equation}
Carrying out the matrix multiplication and extracting $T_{22}$, the
transmission amplitude is
\begin{equation}
 t = \frac{1}{T_{22}}
 = \frac{4n\,e^{-ink_0 d}}{(1+n)^2 - (1-n)^2\,e^{2ink_0 d}},
 \label{eq:fp-transmission-amplitude}
\end{equation}
and the transmission intensity is the Airy function
\begin{equation}
 |t|^2 = \frac{1}{1 + F\sin^2(nk_0 d)},
 \qquad F = \frac{4R}{(1-R)^2},
 \label{eq:airy-from-transfer}
\end{equation}
where $R = |(n-1)/(n+1)|^2$ is the single-interface power reflectance
and $F$ is the coefficient of finesse. This recovers the result
derived from physical arguments in Chapter~\ref{ch:boundaries} and
provides a systematic framework for extending the calculation to
multilayer stacks.

\paragraph{Resonance condition and QNM connection.}
The resonances of the Fabry--P\'erot cavity correspond to the poles of
the transmission amplitude, i.e.\ the zeros of $T_{22}(\omega)$.
From~\eqref{eq:fp-transmission-amplitude}, the pole condition is
$(1+n)^2 = (1-n)^2\,e^{2ink_0 d}$, which gives
\begin{equation}
 \tilde{\omega}_m = \frac{c}{nd}\left(m\pi - i\ln\frac{1+n}{|1-n|}\right),
 \qquad m = 0, \pm 1, \pm 2, \ldots
 \label{eq:fp-poles-transfer}
\end{equation}
The real part gives the resonance frequencies (equally spaced by the
free spectral range $\pi c/(nd)$); the imaginary part gives the decay
rate, which equals the QNM linewidth $\gamma = c\,\ln|(1+n)/(1-n)|/(nd)$
derived in Chapter~\ref{ch:qnm}. The transfer-matrix approach thus
provides a systematic route from multilayer geometry to QNM frequencies,
applicable to arbitrary one-dimensional photonic structures including
those with gain, loss, and dispersion.
